% =========================================================================
% 1. DOCUMENT CLASS & BASIC LAYOUT CONFIGURATION
% =========================================================================
\usepackage[margin=1.0in]{geometry}
\setlength{\marginparwidth}{2cm} % Explicitly sets width for margin notes (todonotes)
\usepackage{multicol}            % Standard multi-column formatting
\usepackage{adjustbox}           % Sizing utilities, cropping, and adjustments

\numberwithin{equation}{section} 
%\AtBeginDocument{\def\MR#1{}}   % Suppresses Mathematical Reviews (MR) indexing tags from citations

% =========================================================================
% 2. CORE TYPOGRAPHY & DOCUMENT FONT ENGINES
% =========================================================================
\usepackage[utf8]{inputenc}
\usepackage[T1]{fontenc}
%\usepackage{CormorantGaramond}
\usepackage{lmodern}     % Latin Modern typography engine
\usepackage{libertinus}  % Libertinus font family suite
\usepackage{microtype}   % Automated spacing, justification, and micro-kerning adjustments
\usepackage{xspace}      % Dynamic tracking spacing for inline macros

% =========================================================================
% 3. COMPREHENSIVE MATHEMATICAL TOOLKITS
% =========================================================================
\usepackage{amsmath}
\usepackage{amssymb}
\usepackage{amsfonts}
\usepackage{amsthm}
\usepackage{faktor}
\usepackage{mathtools}   % Adds extended utilities like \coloneqq
\usepackage{mathrsfs}    % Script math formatting
%\usepackage{mathptmx}    % Times-like font mapping
\usepackage{newtxtext}
\usepackage{newtxmath}   % Math layout fonts mirroring standard text types
\usepackage{mathdots}    % Adds variant diagonal math dots
\usepackage{stackrel}    % Multi-level positional stacked symbols
\usepackage{bbm}         % Extra variant blackboard bold structures (e.g. \mathbbm{1})
\usepackage{latexsym}    % Legacy math symbol parameters
\usepackage{cases}
\usepackage{csquotes}
\usepackage{etoolbox}
\usepackage{mathrsfs}
\usepackage{anyfontsize}
% =========================================================================
% 4. INTERACTIVE REFERENCES & DOCUMENT UTILITIES
% =========================================================================
\usepackage{enumitem}    % Custom layout for lists
\usepackage{graphicx}    % Graphics toolkit for loading imagery or utilizing \scalebox
\usepackage{array}       % Advanced modifications for column dimensions
\usepackage{comment}     % Enables block execution nesting via \begin{comment} ... \end{comment}
\usepackage{xcolor}      % Explicit customized palette rendering tools

% Color Mapping
\definecolor{cherryred}{HTML}{D2042D}

% Hyperlink Engine

%\usepackage[colorlinks=true, linkcolor=cherryred, citecolor=cherryred, urlcolor=blue]{hyperref}
% Clever Cross-Referencing
\usepackage[nameinlink,capitalize]{cleveref}

% =========================================================================
% 5. ADVANCED TIKZ VECTOR GRAPHICS & COMMUTATIVE DIAGRAM ENVIRONMENTS
% =========================================================================
\usepackage{tikz}
\usepackage{tikz-3dplot}
\usepackage{quiver}      % Custom category framework engine matching online quiver utilities
\usepackage{tikz-cd}

% TikZ Libraries Import
\usetikzlibrary{cd}
\usetikzlibrary{nfold}
\usetikzlibrary{math}
\usetikzlibrary{fit}
\usetikzlibrary{calc}
\usetikzlibrary{arrows}
\usetikzlibrary{patterns}
\usetikzlibrary{positioning}
\usetikzlibrary{intersections}
\usetikzlibrary{matrix}
\usetikzlibrary{decorations.pathmorphing}  % Squiggly arrows
\usetikzlibrary{decorations.pathreplacing} % Curly brackets
\usetikzlibrary{decorations.markings}      % Dual path marks and arrow parameters

% Custom Higher-Category Arrow Configurations (2-cells, 3-cells, 4-cells)
\tikzset{
    Rightarrow/.style={double equal sign distance,>={Implies},->},
    triple/.style={-,preaction={draw,Rightarrow}},
    quadruple/.style={preaction={draw,Rightarrow,shorten >=0pt},shorten >=1pt,-,double,double distance=0.2pt}
}

% Adjoint Sign Positioning Tool (\bot sign generation inside paths without a visible line)
\tikzset{
    symbol/.style={
        draw=none,
        every to/.append style={
            edge node={node [sloped, allow upside down, auto=false]{$#1$}}}
    }
}

\newcommand{\adjunction}[4]{%
  \begin{tikzcd}[
    ampersand replacement=\&,
    baseline=(current bounding box.center),
    column sep=5em
  ]
  #1
    \arrow[r, shift left=1.2ex, "#3"]
    \arrow[r, phantom, "\bot" description]
  \&
  #2
    \arrow[l, shift left=1.2ex, "#4"]
  \end{tikzcd}%
}

% Scaled Arrow Configurations for 2-cells
\tikzset{
    scalearrow/.style n args={3}{
        decoration={
            markings,
            mark=at position (1-#1)/2*\pgfdecoratedpathlength with {\coordinate (#2);},
            mark=at position (1+#1)/2*\pgfdecoratedpathlength with {\coordinate (#3);},
        },
        postaction=decorate,
    }
}

% =========================================================================
% =========================================================================
% 6. DYNAMIC LOOP COUNTER (SAFE AUTOMATIC FONT VARIABLE GENERATION)
% =========================================================================
\makeatletter
% Automatically loops over all 26 Latin capital letters to generate macros safely.
% Only creates \b[A-Z], \c[A-Z], and \bf[A-Z] if they don't already exist.
\newcommand{\@bbify}[1]{
  \ifcsname b#1\endcsname
    \message{WARNING: Skipping b#1 to prevent infinite loop conflict!}
  \else
    \expandafter\edef\csname b#1\endcsname{\noexpand\ensuremath{\noexpand\mathbb #1}\noexpand\xspace}
  \fi}

\newcommand{\@calify}[1]{
  \ifcsname c#1\endcsname
    \message{WARNING: Skipping c#1 to prevent infinite loop conflict!}
  \else
    \expandafter\edef\csname c#1\endcsname{\noexpand\ensuremath{\noexpand\mathcal #1}\noexpand\xspace}
  \fi}

\newcommand{\@bfify}[1]{
  \ifcsname bf#1\endcsname
    \message{WARNING: Skipping bf#1 to prevent infinite loop conflict!}
  \else
    \expandafter\edef\csname bf#1\endcsname{\noexpand\ensuremath{\noexpand\mathbf #1}\noexpand\xspace}
  \fi}

\newcounter{@letter}\stepcounter{@letter}
\loop\@bbify{\Alph{@letter}}\@calify{\Alph{@letter}}\@bfify{\Alph{@letter}}
\ifnum\the@letter<26\stepcounter{@letter}\repeat
\makeatother
% =========================================================================
% =========================================================================
% 7. UNIFIED THEOREM ENVIRONMENTS & COUNTER RESET BLOCKS
% =========================================================================
% --- Theorem Style: Plain (Italics) ---
\theoremstyle{plain}
\newtheorem{theorem}{Theorem}[section]
\newtheorem{proposition}[theorem]{Proposition}
\newtheorem{lemma}[theorem]{Lemma}
\newtheorem{corollary}[theorem]{Corollary}

% Safely linking alternate shorthand forms to share the master 'theorem' counter sequence
\newtheorem{thm}[theorem]{Theorem} 
\newtheorem{cor}[theorem]{Corollary}
\newtheorem{lem}[theorem]{Lemma}
\newtheorem{prop}[theorem]{Proposition}
\newtheorem{claim}[theorem]{Claim}
\newtheorem{thrm}[theorem]{Theorem}
\newtheorem{conjecture}[theorem]{Conjecture}
\newtheorem{wconjecture}[theorem]{Wild Conjecture}

% Alpha-Numbered Theorems (e.g., Theorem A, Theorem B)
\newtheorem{thmalph}{Theorem}
\renewcommand*{\thethmalph}{\Alph{thmalph}}
\newtheorem{coralph}{Corollary}
\renewcommand*{\thecoralph}{\Alph{coralph}}

% Unnumbered Plain Statements
\newtheorem*{unnumberedtheorem}{Theorem}
\newtheorem*{unnumberedwannabetheorem}{Wannabe-Theorem}
\newtheorem*{unnumberedcorollary}{Corollary}
\newtheorem*{unnumberedlemma}{Lemma}
\newtheorem*{unnumberedproposition}{Proposition}
\newtheorem*{unnumberedclaim}{Claim}

% --- Theorem Style: Definition (Upright Body text, Bold Header) ---
\theoremstyle{definition}
\newtheorem{definition}[theorem]{Definition}
\newtheorem{construction}[theorem]{Construction}
\newtheorem{notation}[theorem]{Notation}

\newtheorem{defn}[theorem]{Definition}
\newtheorem{const}[theorem]{Construction}
\newtheorem{notn}[theorem]{Notation}
\newtheorem{convention}[theorem]{Convention}
\newtheorem{question}[theorem]{Question}
\newtheorem{rmk}[theorem]{Remark}
\newtheorem{ex}[theorem]{Example}

\newtheorem*{unnumberedtheoremdefinition}{Definition}

% --- Theorem Style: Remark (Upright Body text, Italic Header) ---
\theoremstyle{remark}
\newtheorem{remark}[theorem]{Remark}
\newtheorem{hypothesis}[theorem]{Hypothesis}
\newtheorem{recall}[theorem]{Recall}
\newtheorem{nonex}[theorem]{Non-Example}
\newtheorem{observation}[theorem]{Observation}
\newtheorem{warning}[theorem]{Warning}
\newtheorem{digression}[theorem]{Digression}
\newtheorem{answ}[theorem]{Answer}
\newtheorem{crzcond}[theorem]{Crazy Condition}
\newtheorem{aside}[theorem]{Aside}

\newtheorem*{unnumberedtheoremremark}{Remark}
\newtheorem*{unnumberedtheoremexample}{Example}

% Master Equations Counter Configurations (Tying equations to the theorem counter)
\makeatletter
\let\c@equation\c@theorem
\numberwithin{equation}{section}
\makeatother

% =========================================================================
% 8. INTELLIGENT CROSS-REFERENCING SETTINGS (cleveref)
% =========================================================================
\crefname{lem}{Lemma}{Lemmas}
\crefname{thm}{Theorem}{Theorems}
\crefname{defn}{Definition}{Definitions}
\crefname{notn}{Notation}{Notations}
\crefname{note}{Note}{Note}
\crefname{const}{Construction}{Constructions}
\crefname{prop}{Proposition}{Propositions}
\crefname{rmk}{Remark}{Remarks}
\crefname{goal}{Goal}{Goals}
\crefname{cor}{Corollary}{Corollaries}
\crefname{equation}{Display}{Displays}
\crefname{ex}{Example}{Examples}
\crefname{nonex}{Non-Example}{Non-Examples}
\crefname{thmalph}{Theorem}{Theorems}
\crefname{answ}{Answer}{Answers}
\crefname{question}{Question}{Questions}

% =========================================================================
% 9. CATEGORY THEORY MACROS & ABSTRACT STRUCTURES
% =========================================================================
% Standard Bold Namespaces
\newcommand{\Cat}{\mathbf{Cat}}
\newcommand{\Strat}{\mathbf{Strat}}
\newcommand{\qcat}{Q\mathbf{Cat}}
\newcommand{\sqcat}{sQ\mathbf{Cat}}
\newcommand{\sStrat}{s\mathbf{Strat}}
\newcommand{\Sat}{\mathbf{Sat}}
\newcommand{\mset}{\mathbf{Set}^{+}_{\Delta}}
\newcommand{\sset}{\mathbf{sSet}}
\newcommand{\ssset}{\mathbf{ssSet}}
\newcommand{\rmap}{\mathbf{R}\mathrm{Map}}
\newcommand{\reedyslice}{\bigl((\sset_{(\infty,1)})^{\Delta^{\mathrm{op}}}_{/\bisimFin}\bigr)}
\newcommand{\genreedyslice}{\mathbf M^{\Delta^{\mathrm{op}}}_{/\mathbb B}}
\newcommand{\cocart}{\bigl(\sset^+_{/(\simFin)^\sharp}\bigr)_{\mathrm{coCart}}}
\newcommand{\cocartsp}{\bigl(\sset^+_{/(\simFin)^\sharp}\bigr)_{\mathrm{coCart}}^{\mathrm{sp}}}
\newcommand{\genldblfib}{\bigl((\Minfn)^{\Delta^{\mathrm{op}}}_{/\bisimFin}\bigr)_{\mathrm{LDbl}}}
\newcommand{\genldblfibsp}{\bigl((\Minfn)^{\Delta^{\mathrm{op}}}_{/\bisimFin}\bigr)_{\mathrm{LDbl}}^{\mathrm{sp}}}
\newcommand{\ldblfib}{\bigl((\sset_{(\infty,1)})^{\Delta^{\mathrm{op}}}_{/\bisimFin}\bigr)_{\mathrm{LDbl}}}
\newcommand{\ldblfibsp}{\bigl((\sset_{(\infty,1)})^{\Delta^{\mathrm{op}}}_{/\bisimFin}\bigr)_{\mathrm{LDbl}}^{\mathrm{sp}}}

% Script Category Namespaces
\newcommand{\cat}{\cC\!\mathit{at}}
\newcommand{\set}{\cS\!\mathit{et}}
\newcommand{\ch}{\mathcal C\mathit{h}}
\newcommand{\ach}{\mathit{a}\ch}
\newcommand{\adch}{\mathit{ad}\mathcal C\mathit{h}}
\newcommand{\msset}{m\sset}
\newcommand{\twocat}{2\cat}
\newcommand{\omegacat}{\omega\cat}
\newcommand{\vcat}[1]{\cat_{#1}}
\newcommand{\Set}{\mathbf{Set}}
\newcommand{\sSet}{\mathbf{sSet}}
%\newcommand{\Cat}{\mathbf{Cat}}
\newcommand{\TwoCat}{\mathbf{2Cat}}
\newcommand{\Fin}{\mathrm{Fin}_*}
\newcommand{\Dop}{\mathbf{\Delta}^{\mathrm{op}}}
\newcommand{\CMon}{\mathrm{CMon}}
\newcommand{\bbE}{\mathbb{E}}
\newcommand{\Htwo}{\mathbb{H}_2}
\newcommand{\Nh}{N^{\mathrm{h}}}
\newcommand{\Iso}{\mathrm{Iso}}
\newcommand{\Hor}{\mathrm{Hor}}
\newcommand{\Pres}{\mathrm{Pres}}
\newcommand{\Ex}{\mathrm{Ex}}
\newcommand{\Seg}{\mathrm{Seg}}
\newcommand{\PIso}{\mathrm{P}^{\cong}}
\newcommand{\PEq}{\mathrm{P}^{\simeq}}
\newcommand{\LL}{\mathbf{L}}
\newcommand{\RR}{\mathbf{R}}
\newcommand{\nat}{\mathrm{nat}}
\newcommand{\Lk}{\mathrm{L}}
\newcommand{\llp}{\mathrm{llp}}
\newcommand{\rlp}{\mathrm{rlp}}
\newcommand{\pt}[1]{\langle #1\rangle}
\newcommand{\hto}{\rightsquigarrow}
\newcommand{\hc}{\odot}
\newcommand{\hid}[1]{1^{\mathrm{h}}_{#1}}


% Specialized Script Elements
\newcommand{\C}{\mathcal{C}}
\newcommand{\B}{\mathcal{B}}
\newcommand{\E}{\mathcal{E}}
\newcommand{\M}{\mathcal{M}}
\newcommand{\mrt}{\mathfrak M}
\newcommand{\G}{\mathcal G}
\newcommand{\I}{\mathcal I}
\newcommand{\RMap}{\mathbf{R}\mathrm{Map}}
\newcommand{\one}{\mathsf 1}

% Structural Functor Mapping Shorthands
\newcommand{\F}{\mathcal{F}_*}
\newcommand{\simFin}{\mathcal{F}\mathrm{in}_*}
\newcommand{\bisimFin}{\mathbb{F}\mathrm{in}_*}
\newcommand{\psh}[1]{\set^{#1^{\op}}}
\newcommand{\spsh}[1]{\sset^{#1^{\op}}}
\newcommand{\comp}{\ast}
\newcommand{\Fun}{\mathbf{Fun}} 
\newcommand{\ihom}[2]{#2^{#1}}
\newcommand{\vect}[1]{\vec{#1}}
\newcommand{\cof}{\hookrightarrow}
\newcommand{\fib}{\twoheadrightarrow}
\newcommand{\trivcof}{\overset{\sim}{\hookrightarrow}}
\newcommand{\trivfib}{\overset{\sim}{\twoheadrightarrow}}

% =========================================================================
% 10. MATH OPERATORS & SPECIALIZED TRUNCATION FUNCTIONS
% =========================================================================
% Standard Categories Definitions & Mappings
\DeclareMathOperator{\Ho}{Ho}
\DeclareMathOperator{\map}{map}
\DeclareMathOperator{\Hom}{Hom}
\DeclareMathOperator{\Map}{Map} 
\DeclareMathOperator{\Mor}{Mor}
\DeclareMathOperator{\Ob}{Ob}
\DeclareMathOperator{\Cocart}{Cocart}


% General Math Maps & Substructures
\DeclareMathOperator{\incl}{incl}
\DeclareMathOperator{\tr}{tr}
\DeclareMathOperator{\core}{core}
\DeclareMathOperator{\colim}{colim}
\DeclareMathOperator{\holim}{holim}
\DeclareMathOperator{\coeq}{coeq}
\DeclareMathOperator{\id}{id}
\DeclareMathOperator{\im}{im}
\DeclareMathOperator{\op}{op}
\DeclareMathOperator{\ev}{ev}
\DeclareMathOperator{\pr}{pr}
\DeclareMathOperator{\N}{N}
\DeclareMathOperator{\diag}{diag}
\DeclareMathOperator{\Sp}{Sp}

% Custom Identity Spacings
\newcommand{\ide}{\mathrm{id}\mathord{}}
\newcommand{\idf}{\mathrm{id}\,}
\newcommand{\idg}{\mathrm{id}\hspace{0.1111em}}

% Simplicial Combinatorics Operators
\newcommand{\aamalg}[1]{\underset{#1}{{\amalg}}} 
\newcommand{\eqDelta}{\Delta[3]_{\mathrm{eq}}}
\newcommand{\thin}[2]{t#1_{#2}}
\newcommand{\maybethin}[2]{(t)#1_{#2}}
\newcommand{\thn}{\mathrm{th}}

% Specialized Indexing Shapes
\newcommand{\Osim}[1]{\cO_2^\sim[{#1}]}
\newcommand{\DelQuot}{\overline{\Delta}}
\newcommand{\AdjEq}{\mathrm{AdjEq}}
\newcommand{\Otilde}{\widetilde{\cO}}
\newcommand{\ortl}[1]{\cO[#1]}

% Truncation Operations
\newcommand{\chaos}{\mathrm{ch}}
\newcommand{\roughtrunc}[1]{\tau_{\leq #1}^{\mathrm{b}}}
\newcommand{\inttrunc}[1]{\tau_{\leq #1}^{\mathrm{i}}}


% ===== Macros for "Symmetric monoidal (infty,n)-categories" =====
% These names avoid changing the shared meanings of \M, \G, \I and \LL.
\newcommand{\Minfn}{\mathbf M_{(\infty,n)}}
\newcommand{\HGen}{\mathsf G}
\newcommand{\Stair}{\mathcal{S}\mathrm{tair}}
\newcommand{\Int}{\mathbb{I}\mathrm{nt}}
\newcommand{\Diag}{\mathcal{D}\mathrm{iag}}
\newcommand{\diagset}{\mathrm{diag}}
% Full environment names used by the manuscript.
\crefname{construction}{Construction}{Constructions}
\crefname{notation}{Notation}{Notations}
\crefname{convention}{Convention}{Conventions}
\crefname{definition}{Definition}{Definitions}
\crefname{remark}{Remark}{Remarks}
\crefname{hypothesis}{Hypothesis}{Hypotheses}
% ===== End of paper-specific macros =====
